% main.tex
\documentclass[a4paper,12pt]{article}

% ---------------- MARGINS ----------------
\usepackage[left=2.36cm, right=2.33cm, top=1.75cm, bottom=1.27cm]{geometry}

% ---------------- PACKAGES ----------------
\usepackage{fancyhdr}
\usepackage{graphicx}
\usepackage{multirow}
\usepackage{tikz}
\usepackage{everypage}
\usepackage{setspace}
\usepackage{multicol}
\usepackage{algorithm}
\usepackage{algorithmic}
\usepackage{subcaption}
\usepackage{amsmath}
\usepackage{booktabs}
\usepackage{hyperref}
\usepackage{cite}
\usepackage{caption}
\usepackage{xcolor}
\usepackage{times}
\usepackage{tocloft}
\usepackage{ragged2e}
\usepackage{float}
\renewcommand{\cftdot}{}
\justifying
\usepackage{chngcntr}
\counterwithin{figure}{section}
\usetikzlibrary{positioning}
\captionsetup{font=small}

% ---------------- REMOVE RED BOXES ----------------
\hypersetup{
    colorlinks=false,
    pdfborder={0 0 0}
}

% ---------------- TOC FORMATTING ----------------
\renewcommand{\cftsecfont}{\bfseries}
\renewcommand{\cftsecpagefont}{\bfseries}
\renewcommand{\cftsubsecfont}{\normalfont}
\renewcommand{\cftsubsecpagefont}{\normalfont}
\setlength{\cftsecindent}{0em}
\setlength{\cftsecnumwidth}{2.5em}
\setlength{\cftsubsecindent}{2.5em}
\setlength{\cftsubsecnumwidth}{3em}

% ---------------- PAGE BORDER ----------------
\AddEverypageHook{%
    \begin{tikzpicture}[remember picture, overlay]
        \draw[line width=1pt]
            ([xshift=10mm, yshift=15mm] current page.south west)
            rectangle
            ([xshift=-10mm, yshift=-15mm] current page.north east);
    \end{tikzpicture}%
}

% ---------------- SECTION HEADING SIZES (match docx 12pt body) ----------------
\usepackage{titlesec}
\titleformat{\section}{\normalfont\fontsize{12}{14}\bfseries}{\thesection}{1em}{}
\titleformat{\subsection}{\normalfont\fontsize{12}{14}\bfseries}{\thesubsection}{1em}{}
\titleformat{\subsubsection}{\normalfont\fontsize{12}{14}\bfseries}{\thesubsubsection}{1em}{}
\titlespacing*{\section}{0pt}{12pt}{6pt}
\titlespacing*{\subsection}{0pt}{10pt}{4pt}
\titlespacing*{\subsubsection}{0pt}{8pt}{4pt}

% ---------------- FONT & SPACING ----------------
\onehalfspacing

% ---------------- HEADER HEIGHT FIX ----------------
\setlength{\headheight}{15pt}

% ---------------- PLAIN STYLE ----------------
\fancypagestyle{plain}{
    \fancyhf{}
    \fancyfoot[L]{Dept. of CST, DSU}
    \fancyfoot[R]{\thepage}
    \renewcommand{\headrulewidth}{0pt}
    \renewcommand{\footrulewidth}{0pt}
}

% ---------------- MAIN PAGE STYLE ----------------
\pagestyle{fancy}
\fancyhf{}
\fancyhead[L]{\small\textbf{Noise-Aware Pipeline for Indian Bird Sound Classification Using BirdNET Embeddings}}
\fancyfoot[L]{Dept. of CST, DSU}
\fancyfoot[R]{\thepage}
\renewcommand{\headrulewidth}{0pt}
\renewcommand{\footrulewidth}{0pt}

\begin{document}
\pagenumbering{gobble}

% -------------------------
% COVER PAGE
% -------------------------
\thispagestyle{empty}

\begin{center}

\begin{minipage}{0.4\textwidth}
\centering
\includegraphics[width=0.8\linewidth]{logo_d.jpg}
\end{minipage}
\hfill
\begin{minipage}{0.4\textwidth}
\centering
\includegraphics[width=0.8\linewidth]{logo_s.png}
\end{minipage}

\vspace{0.8cm}

{\fontsize{32}{36}\selectfont\textbf{Dayananda Sagar University}} \\[4pt]
{\fontsize{24}{28}\selectfont\textbf{School of Engineering}} \\[6pt]
{\fontsize{12}{14}\selectfont\textbf{Devarakaggalahalli, Harohalli, Kanakapura Road,\\
Bangalore South Dt., Bengaluru -- 562112}}

\vspace{0.8cm}

{\fontsize{26}{30}\selectfont\textbf{Department of}}\\[2pt]
{\fontsize{26}{30}\selectfont\textbf{Computer Science \& Technology}}

\vspace{0.8cm}

{\fontsize{15}{18}\selectfont\textbf{ACADEMIC YEAR 2025--26}}

\vspace{0.8cm}

{\fontsize{20}{24}\selectfont\textbf{Minor Project (23CT3613)}} \\[4pt]
{\fontsize{20}{24}\selectfont\textbf{Report on}}

\vspace{0.6cm}

{\fontsize{20}{24}\selectfont\textbf{
Noise-Aware Pipeline for Indian Bird Sound Classification\\
Using BirdNET Embeddings and Learned Gating
}}

\vspace{1cm}

{\fontsize{14}{16}\selectfont By}

\vspace{0.4cm}

\begin{tabular}{ll}
{\fontsize{15}{18}\selectfont\textbf{Samarth D Kolur}} & {\fontsize{15}{18}\selectfont\textbf{ENG23CT0019}} \\[4pt]
{\fontsize{15}{18}\selectfont\textbf{Manoj C}}         & {\fontsize{15}{18}\selectfont\textbf{ENG23CT0031}} \\[4pt]
{\fontsize{15}{18}\selectfont\textbf{Mithun Krishnan}} & {\fontsize{15}{18}\selectfont\textbf{ENG23CT0034}} \\[4pt]
{\fontsize{15}{18}\selectfont\textbf{Yashaswini R}}    & {\fontsize{15}{18}\selectfont\textbf{ENG23CT0044}} \\
\end{tabular}

\vspace{1cm}

{\fontsize{15}{18}\selectfont\textbf{Under the supervision of}} \\[6pt]
{\fontsize{14}{16}\selectfont\textbf{Dr. Ramandeep Kaur}} \\[2pt]
{\fontsize{14}{16}\selectfont\textbf{Assistant Professor,}} \\
{\fontsize{14}{16}\selectfont\textbf{Department of Computer Science and Technology}} \\
{\fontsize{14}{16}\selectfont\textbf{Dayananda Sagar University}}

\vspace{1cm}

{\fontsize{14}{16}\selectfont\textbf{DEPARTMENT OF COMPUTER SCIENCE \& TECHNOLOGY}} \\
{\fontsize{14}{16}\selectfont\textbf{DAYANANDA SAGAR UNIVERSITY}}

\end{center}

% -------------------------
% CERTIFICATE PAGE
% -------------------------
\newpage
\thispagestyle{empty}
\vspace{2cm}
\begin{center}
\includegraphics[width=0.4\textwidth]{logo_d.jpg}
\hspace{70pt}
\includegraphics[width=0.4\textwidth]{logo_s.png}

\vspace{1cm}
{\fontsize{36}{40}\selectfont\textbf{Dayananda Sagar University}} \\[4pt]
{\fontsize{24}{28}\selectfont\textbf{School of Engineering}} \\[4pt]
{\fontsize{12}{14}\selectfont\textbf{Devarakaggalahalli, Harohalli, Kanakapura Road, Bangalore South Dt., Bengaluru -- 562112}}

\vspace{0.8cm}
{\fontsize{20}{24}\selectfont\textbf{Department of Computer Science \& Technology}}

\vspace{0.8cm}
{\fontsize{18}{20}\selectfont\textbf{CERTIFICATE}}
\end{center}

\vspace{1cm}

\justify
This is to certify that the work titled
``\textbf{Noise-Aware Pipeline for Indian Bird Sound Classification Using BirdNET Embeddings and Learned Gating}''
is carried out by
\textbf{Samarth D Kolur -- ENG23CT0019},
\textbf{Manoj C -- ENG23CT0031},
\textbf{Mithun Krishnan -- ENG23CT0034},
and \textbf{Yashaswini R -- ENG23CT0044},
bonafide students of Bachelor of Technology in Computer Science and Technology at the School of Engineering,
Dayananda Sagar University, Bangalore, in partial fulfillment for the award of degree in Bachelor
of Technology in Computer Science and Technology during the academic year 2025--2026.

\vspace{2cm}

\noindent
\begin{tabular*}{\textwidth}{@{\extracolsep{\fill}} c c}
\textbf{Dr. Ramandeep Kaur} & \textbf{Dr. Shahina M Parveen} \\[6pt]
\begin{tabular}{c}
Assistant Professor,\\
Dept. of CST,\\
School of Engineering,\\
Dayananda Sagar University.
\end{tabular}
&
\begin{tabular}{c}
Chairperson, CST,\\
Dept. of CST,\\
School of Engineering,\\
Dayananda Sagar University.
\end{tabular}
\end{tabular*}

\vspace{1.5cm}

\noindent \textbf{Date:}

% -------------------------
% DECLARATION PAGE
% -------------------------
\newpage
\thispagestyle{plain}

\vspace{1.5cm}

\begin{center}
\includegraphics[width=0.4\textwidth]{logo_d.jpg}
\hspace{70pt}
\includegraphics[width=0.4\textwidth]{logo_s.png}

\vspace{0.8cm}

{\fontsize{36}{40}\selectfont\textbf{Dayananda Sagar University}} \\[4pt]
{\fontsize{24}{28}\selectfont\textbf{School of Engineering}} \\[4pt]
{\fontsize{12}{14}\selectfont\textbf{Devarakaggalahalli, Harohalli, Kanakapura Road, Bangalore South Dt., Bengaluru -- 562112}}

\vspace{0.8cm}
{\fontsize{20}{24}\selectfont\textbf{Department of Computer Science \& Technology}}

\vspace{0.8cm}
{\fontsize{14}{16}\selectfont\textbf{DECLARATION}}
\end{center}

\vspace{0.8cm}

\justify
We, \textbf{Samarth D Kolur -- ENG23CT0019},
\textbf{Manoj C -- ENG23CT0031},
\textbf{Mithun Krishnan -- ENG23CT0034},
and \textbf{Yashaswini R -- ENG23CT0044},
are students of the \textbf{sixth semester} B.Tech in Computer Science and Technology at the
School of Engineering, Dayananda Sagar University. We hereby declare that the minor project titled
``\textbf{Noise-Aware Pipeline for Indian Bird Sound Classification Using BirdNET Embeddings and Learned Gating}''
has been carried out by us and submitted in partial fulfillment for the award of degree
in Bachelor of Technology in Computer Science and Technology during the academic year 2025--2026.

\vspace{0.8cm}

\begin{tabular}{p{0.08\textwidth} p{0.5\textwidth} p{0.32\textwidth}}
\textbf{Sl.No.} & \textbf{Name and USN} & \textbf{Student Signature} \\[8pt]
1 & Samarth D Kolur \newline ENG23CT0019 & \\[12pt]
2 & Manoj C \newline ENG23CT0031 & \\[12pt]
3 & Mithun Krishnan \newline ENG23CT0034 & \\[12pt]
4 & Yashaswini R \newline ENG23CT0044 & \\
\end{tabular}

\vspace{1.5cm}

\noindent
\hspace{0.5cm} Place: DSU, Bangalore \hfill Date:

% -------------------------
% ACKNOWLEDGEMENT
% -------------------------
\newpage
\thispagestyle{plain}

\begin{center}
{\fontsize{14}{16}\bfseries ACKNOWLEDGEMENT}
\end{center}

\vspace{0.5cm}

\justify
\noindent
It is a great pleasure for us to acknowledge the assistance and support of many individuals who have been
responsible for the successful completion of this \textbf{Noise-Aware Pipeline for Indian Bird Sound
Classification using BirdNET Embeddings and Learned Gating}. First, we take this opportunity to express
our sincere gratitude to the School of Engineering and Technology, Dayananda Sagar University for
providing us with a great opportunity to pursue our bachelor's degree in this institution.

\vspace{0.3cm}

\noindent
We would like to thank \textbf{Dr. Udaya Kumar Reddy K R, Dean, School of Engineering and Technology,
Dayananda Sagar University} for providing this opportunity.

\vspace{0.3cm}

\noindent
It is an immense pleasure to express our sincere thanks to \textbf{Dr. M. Shahina Parveen, Chairperson,
Department of Computer Science and Technology, Dayananda Sagar University}, for providing the right
academic guidance and motivating us during the course.

\vspace{0.3cm}

\noindent
We express our sincere thanks to our guide, \textbf{Dr. Ramandeep Kaur}, Assistant Professor,
Department of Computer Science and Technology, for her valuable time, academic guidance, and constant
support throughout this project.

% -------------------------
% ABSTRACT
% -------------------------
\newpage
\thispagestyle{plain}

\begin{center}
{\fontsize{14}{16}\bfseries ABSTRACT}
\end{center}

\addcontentsline{toc}{section}{Abstract}

\vspace{0.5cm}

\justify
\noindent
The increasing use of passive acoustic monitoring has generated large volumes of environmental audio
data, making automated bird sound classification both essential and challenging. Real-world recordings
often contain significant background noise such as wind, insects, and human activity, which reduces
the accuracy of traditional classification approaches that rely on direct model predictions without
noise handling.

\vspace{0.3cm}

\noindent
This project presents a \textbf{noise-aware bird sound classification pipeline} designed to improve
reliability in real-world conditions. The system processes raw field recordings by standardizing
audio, segmenting it into fixed durations, and removing low-energy segments using RMS-based
filtering. A dedicated \textbf{Noise Segregation V2} module further analyzes acoustic features
such as zero-crossing rate, spectral flatness, spectral centroid, and periodicity across subframes,
using a majority voting mechanism to distinguish between bird-like signals and background noise.

\vspace{0.3cm}

\noindent
For feature representation, the system utilizes \textbf{BirdNET embeddings}, extracting
high-dimensional audio features from a pre-trained model. These embeddings are then passed to a
compact multilayer perceptron (MLP) classifier trained using focal loss for improved performance
on imbalanced data. An optional \textbf{autoencoder-based validation mechanism} is incorporated
to detect out-of-distribution inputs by analyzing reconstruction error, thereby enhancing
robustness.

\vspace{0.3cm}

\noindent
During inference, a three-level decision strategy is applied to classify segments into bird,
noise, or uncertain categories, with fixed thresholds $\tau_{\mathrm{low}}=0.3$ and
$\tau_{\mathrm{high}}=0.7$. The system also includes an \textbf{active learning component}
that identifies difficult or ambiguous samples for further manual review, enabling continuous
improvement of the model.

\vspace{0.3cm}

\noindent
Experimental results on the iBC53 dataset demonstrate significant performance improvements
across 7,193 test segments, with accuracy increasing from 0.570 to 0.992 and F1-score from
0.701 to 0.995, along with an 88\% reduction in false-positive rate and a 99\% reduction in
false-negative rate relative to the BirdNET confidence baseline.

\vspace{0.5cm}

\noindent
\textbf{Keywords:} Bird Sound Classification, Bioacoustics, Noise Segregation, BirdNET,
Machine Learning, MLP Classifier, Autoencoder, Active Learning, Acoustic Feature Analysis

% -------------------------
% TABLE OF CONTENTS
% -------------------------
\newpage
\tableofcontents

% -------------------------
% LIST OF FIGURES (separate page)
% -------------------------
\newpage
\listoffigures

% -------------------------
% LIST OF TABLES (separate page)
% -------------------------
\newpage
\listoftables

% -------------------------
% PAGE NUMBER START
% -------------------------
\newpage
\pagenumbering{arabic}
\setcounter{page}{1}

% ======================================================
% CHAPTER 1
% ======================================================
\newpage
\thispagestyle{plain}

\begin{flushleft}
{\fontsize{16}{18}\bfseries CHAPTER 1}
\end{flushleft}
\begin{center}
{\fontsize{16}{18}\bfseries INTRODUCTION}
\end{center}

\addcontentsline{toc}{section}{\textbf{CHAPTER 1}}
\vspace{0.5cm}

\setcounter{section}{0}
\renewcommand{\thesection}{1.\arabic{section}}

% ---------------- 1.1 ----------------
\section{Introduction}

\noindent
Acoustic monitoring has become an increasingly important method for studying biodiversity, especially
in environments where direct observation is difficult or impractical. Bird sound classification plays
a crucial role in understanding species behavior, migration patterns, ecological balance, and
environmental change. With the growing availability of passive recording devices, large volumes of
environmental audio data are being collected, making automated analysis both necessary and valuable.

\vspace{0.3cm}

\noindent
However, real-world audio recordings are rarely clean. They often contain various forms of background
noise such as wind, insects, rainfall, traffic, and human activity. These disturbances make it
challenging to accurately identify bird vocalizations, particularly when the sounds are faint or
overlap with other acoustic signals. As a result, even advanced systems may struggle to differentiate
between actual bird sounds and noise.

\vspace{0.3cm}

\noindent
One of the most widely adopted systems for bird sound recognition is BirdNET \cite{BirdNET}, a deep
learning model capable of identifying hundreds of bird species from raw audio recordings. While BirdNET
demonstrates strong species-level performance, it outputs sigmoid-transformed confidence scores that
are not calibrated probabilities. Applying a fixed detection threshold across diverse species and
recording environments results in an elevated rate of both false positives and false negatives,
limiting its reliability for direct deployment in noisy field conditions. This project addresses these
limitations by proposing a structured and noise-aware pipeline that integrates preprocessing, feature
extraction, and machine learning to improve classification accuracy and reliability in practical
scenarios.

% ---------------- 1.2 ----------------
\section{Motivation}

\noindent
The motivation for this project arises from the practical limitations observed in existing bird sound
classification systems. In real-world environments, audio recordings are often unstructured and
contain significant background noise, especially in regions with diverse ecological conditions such as
the Indian subcontinent. Traditional models are generally trained and evaluated under controlled
conditions and are not well-equipped to handle the complexity of field recordings, which leads to
inaccurate predictions and reduced reliability in practical deployment.

\vspace{0.3cm}

\noindent
Another key motivation is the need for automation in large-scale data processing. Manual annotation
of bird sounds is time-consuming and requires significant domain expertise, making it impractical
when working with large volumes of passive acoustic data. There is a clear need for a system that
can automatically and reliably distinguish between meaningful bird vocalizations and background
noise without requiring human intervention at every stage.

\vspace{0.3cm}

\noindent
By combining structured noise filtering, deep feature extraction, and machine learning into a
unified pipeline, the system becomes more robust and adaptable. This enables efficient processing
of large volumes of audio data while maintaining consistent and reliable classification performance
across varying environmental conditions.

% ---------------- 1.3 ----------------
\section{Problem Statement}

\noindent
The primary problem addressed in this project is the accurate classification of bird sounds in
noisy real-world audio recordings. Existing systems often fail to effectively separate bird
vocalizations from background noise, leading to incorrect predictions. Noise segments are frequently
misclassified as bird sounds, inflating false-positive counts, while faint or distant bird calls
are often missed entirely, leading to high false-negative rates. This dual failure mode is a
fundamental limitation of approaches that rely on a single fixed confidence threshold applied
uniformly across all species and environmental conditions.

\vspace{0.3cm}

\noindent
In addition to classification challenges, many systems lack a well-structured pipeline that
integrates all stages of processing, including preprocessing, feature extraction, and
decision-making. This lack of integration leads to inconsistent performance and limits the
practical applicability of such systems in real ecological monitoring tasks.

\vspace{0.3cm}

\noindent
Therefore, the problem can be defined as: designing a reliable and noise-aware system that can
accurately distinguish bird sounds from background noise while maintaining efficiency and
scalability in real-world conditions.

% ---------------- 1.4 ----------------
\section{Objectives}

\noindent
The objectives of this project are designed to develop an efficient and reliable system for bird
sound classification in noisy environments:

\begin{itemize}
    \item To develop a preprocessing framework that standardizes raw audio recordings, segments
          them into fixed durations, and removes low-energy or irrelevant segments.
    \item To design and implement a Noise Segregation V2 mechanism that effectively differentiates
          between bird sounds and background noise using acoustic feature analysis.
    \item To extract meaningful audio features using BirdNET embeddings that improve the
          separability between bird and noise classes.
    \item To build an MLP classification model trained with focal loss that accurately predicts
          whether a given audio segment contains bird sounds or noise.
    \item To incorporate an autoencoder-based validation mechanism to handle uncertain or
          out-of-distribution inputs and improve prediction reliability.
    \item To construct an end-to-end pipeline integrating preprocessing, feature extraction,
          classification, and validation into a structured workflow.
    \item To evaluate the system using accuracy, precision, recall, and F1-score to ensure
          effectiveness in real-world scenarios.
\end{itemize}

% ---------------- 1.5 ----------------
\section{Summary}

\noindent
This chapter introduced the concept of bird sound classification and emphasized its importance in
biodiversity monitoring and ecological studies. It highlighted the challenges associated with
noisy real-world audio data and explained why existing approaches are insufficient for practical
deployment. The motivation for developing a noise-aware and automated system was discussed,
followed by a clear definition of the problem and the specific objectives of the project.

% ======================================================
% CHAPTER 2
% ======================================================
\newpage
\thispagestyle{plain}

\begin{flushleft}
{\fontsize{16}{18}\bfseries CHAPTER 2}
\end{flushleft}
\begin{center}
{\fontsize{16}{18}\bfseries LITERATURE SURVEY}
\end{center}

\addcontentsline{toc}{section}{\textbf{CHAPTER 2}}

\vspace{0.5cm}

\setcounter{section}{0}
\renewcommand{\thesection}{2.\arabic{section}}

% ---------------- 2.1 ----------------
\section{Overview}

\noindent
In recent years, acoustic monitoring has become an important tool for studying biodiversity,
especially in environments where direct observation is difficult. Bird vocalizations provide
valuable insights into species presence, behavior, and ecological changes. With the increasing
use of passive acoustic monitoring systems, large volumes of environmental audio data are being
collected, creating a need for automated systems that can efficiently process and analyze such
data.

\vspace{0.3cm}

\noindent
Real-world audio recordings are often noisy and unstructured. Environmental sounds such as wind,
insects, rainfall, and human activity frequently overlap with bird calls, making accurate
classification challenging. Noise and poor-quality data significantly affect model performance,
leading to incorrect predictions and reduced reliability in practical applications. These
challenges make it necessary to design systems that can effectively handle noise and extract
meaningful information from complex audio signals.

\vspace{0.3cm}

\noindent
To address these issues, modern approaches rely on deep learning techniques such as convolutional
neural networks, which are capable of learning patterns directly from audio spectrograms. These
models have shown strong performance in large-scale bird species recognition tasks, but still
face difficulties in highly noisy environments and depend heavily on well-processed input data.

\vspace{0.3cm}

\noindent
Another important advancement is the use of feature embeddings and transfer learning. Instead
of building models from scratch, embeddings extracted from pre-trained systems can capture
important acoustic characteristics, improving classification accuracy and generalization across
diverse datasets. This approach is especially useful when working with limited or geographically
varied audio collections.

\vspace{0.3cm}

\noindent
Building on these research insights, this project proposes a structured and noise-aware pipeline
for bird sound classification. The system integrates preprocessing, noise segregation, feature
extraction, and classification into a unified workflow, with a validation mechanism to handle
uncertain inputs and improve overall reliability.

% ---------------- 2.2 ----------------
\newpage
\section{Literature Survey}

\noindent
The field of bird sound classification and passive acoustic monitoring has seen significant
advancements through the application of deep learning and transfer learning techniques. Earlier
methods relied on manual observation and expert annotation, which were time-consuming and limited
in scale. The survey below reviews the key works that inform the design of the proposed system.

\vspace{0.3cm}

\noindent
\textbf{Bird Sound Recognition and BirdNET.}
Kahl et al.\ \cite{BirdNET} introduced BirdNET, a large-scale deep learning system for avian
diversity monitoring capable of identifying hundreds of bird species from audio recordings. The
work demonstrates that domain-specific data augmentation and structured preprocessing are critical
for handling noisy and overlapping vocalizations. A subsequent study by Wood and Kahl \cite{WoodKahl}
cautions that BirdNET confidence scores are sigmoid-transformed logits and not calibrated
probabilities; applying a fixed detection threshold across environments routinely inflates both
false-positive and false-negative rates, motivating the need for post-processing pipelines such
as the one proposed here. A broader critical assessment of BirdNET's applications, performance,
and known pitfalls is provided by P\'{e}rez-Granados \cite{PerezGranados}.

\vspace{0.3cm}

\noindent
\textbf{BirdNET Embeddings and Transfer Learning.}
McGinn et al.\ \cite{McGinn} demonstrate that BirdNET embeddings capture meaningful ecological
information beyond species labels and can be used effectively for downstream habitat and
biodiversity analysis. Ghani et al.\ \cite{Ghani} show that global birdsong embeddings derived
from pre-trained models enable superior transfer learning for bioacoustic classification tasks,
significantly outperforming models trained from scratch. Navine et al.\ \cite{Navine} further
demonstrate that embedding-based approaches allow direct estimation of call density in passive
recordings without relying on fixed confidence thresholds.

\vspace{0.3cm}

\noindent
\textbf{Noise Handling and Dataset Preparation.}
Xie et al.\ \cite{Xie} provide a comprehensive review of bioacoustic signal denoising methods,
covering classical signal processing techniques as well as modern deep learning approaches. Their
review confirms that background noise remains a primary factor in degrading model performance and
highlights the importance of structured preprocessing pipelines. The iBC53 dataset used in this
project, comprising passive acoustic recordings from an Indian urban landscape, is described by
Gupta et al.\ \cite{iBC53}, providing ecologically grounded evaluation data for the proposed system.

\vspace{0.3cm}

\noindent
\textbf{Out-of-Distribution Detection.}
Ntalampiras and Potamitis \cite{Ntalampiras} demonstrate reconstruction-error-based methods for
detecting unknown bird species and anomalous acoustic events. Their approach uses autoencoder
networks trained on known-class embeddings, where high reconstruction error signals out-of-distribution
inputs. This principle directly motivates the autoencoder gating mechanism incorporated in the
present system.

\vspace{0.3cm}

\noindent
\textbf{Active Learning for Bioacoustics.}
Kath et al.\ \cite{Kath} show that combining transfer learning with active learning significantly
reduces annotation effort in passive acoustic monitoring, enabling rapid labeling of large audio
corpora with minimal expert involvement. Rauch et al.\ \cite{Rauch} develop deep active learning
strategies specifically for avian bioacoustics, demonstrating that uncertainty-based sample
selection improves classifier calibration over iterative annotation cycles. Dumoulin et al.\
\cite{Dumoulin} further demonstrate that pre-computing embeddings and applying lightweight
active learning classifiers allows rapid model adaptation with minimal expert effort.

\vspace{0.3cm}

\noindent
\textbf{Foundation Models and Loss Functions.}
Broader reviews by Schwinger et al.\ \cite{Schwinger} and van Merri\"{e}nboer et al.\
\cite{Perch2} confirm that pre-trained embedding models represent the current state of the art
for bioacoustic classification, with foundation model architectures demonstrating strong
generalization across taxa and geographies. The focal loss function introduced by Lin et al.\
\cite{FocalLoss} addresses class imbalance during classifier training by down-weighting
easily classified examples, and is adopted in the present system to improve sensitivity
on the minority noise class. The theoretical basis for autoencoder-based dimensionality
reduction and anomaly detection is established by Hinton and Salakhutdinov \cite{Hinton}.

\vspace{0.3cm}

\noindent
Overall, the surveyed works collectively establish that deep learning and embedding-based
transfer learning are highly effective for bird sound classification, while also identifying
persistent challenges around noise robustness and threshold calibration. The proposed system
builds directly on these findings by combining structured noise segregation, BirdNET-based
feature extraction, focal-loss MLP classification, autoencoder gating, and active learning
into a unified and noise-aware pipeline.

% ======================================================
% CHAPTER 3
% ======================================================
\newpage
\thispagestyle{plain}

\begin{flushleft}
{\fontsize{16}{18}\bfseries CHAPTER 3}
\end{flushleft}
\begin{center}
{\fontsize{16}{18}\bfseries SYSTEM DESIGN AND METHODOLOGY}
\end{center}

\addcontentsline{toc}{section}{\textbf{CHAPTER 3}}

\vspace{0.5cm}

\setcounter{section}{0}
\renewcommand{\thesection}{3.\arabic{section}}

% ---------------- 3.1 ----------------
\section{Methodologies}

\noindent
The proposed bird sound classification system follows a structured methodology that integrates
preprocessing, noise segregation, feature extraction, classification, and validation into a
unified pipeline. The system begins with audio preprocessing, where raw recordings are
standardized, segmented into fixed 3-second durations, and filtered to remove silence or
low-energy segments using RMS-based thresholds.

\vspace{0.3cm}

\noindent
To improve data quality, a Noise Segregation V2 module is applied to distinguish between bird
sounds and background noise by analyzing acoustic features such as zero-crossing rate, spectral
flatness, spectral centroid, and insect-band autocorrelation across subframes, using a majority
voting mechanism. This ensures that only meaningful audio data is passed to the embedding stage.

\vspace{0.3cm}

\noindent
Feature extraction is performed using BirdNET V2.4 embeddings, which convert audio signals into
1024-dimensional numerical representations capturing essential acoustic patterns. A multilayer
perceptron (MLP) classifier trained with focal loss is used to predict whether a given segment
contains bird sounds or noise, with early stopping guided by validation F1 score.

\vspace{0.3cm}

\noindent
To further enhance reliability, an optional autoencoder-based validation mechanism is incorporated
to detect out-of-distribution inputs by analyzing reconstruction error. A three-way decision
strategy routes outputs into bird, noise, or uncertain categories, allowing ambiguous cases to
be flagged for human review and iterative retraining.

% ---------------- 3.2 ----------------
\section{Mathematical Model}

\noindent
The proposed system can be represented mathematically as a multi-stage pipeline consisting of
preprocessing, feature extraction, classification, and decision-making.

\subsection{Audio Segmentation}

\noindent
Let the input audio signal be represented as $X = \{x_1, x_2, \ldots, x_n\}$. The signal is
resampled to 48~kHz mono and divided into fixed-length 3-second segments:
\[
    X \rightarrow \{S_1, S_2, \ldots, S_k\}
\]
Any segment whose RMS energy falls below $-40$~dB is discarded before further processing.

\subsection{Noise Segregation V2}

\noindent
Each segment $S_j$ is divided into six non-overlapping 0.5-second subframes. For each subframe
$f_i$, acoustic features are extracted and combined into a weighted noise score:
\[
    \text{Score}(f_i) = w_{\mathrm{zcr}}\,\widehat{\mathrm{ZCR}}
                      + w_{\mathrm{flat}}\,\widehat{\mathrm{Flat}}
                      + w_{\mathrm{cflag}}\,\mathrm{CFlag}
                      + w_{\mathrm{cstd}}\,\hat{\sigma}_c
                      + w_{\mathrm{ins}}\,\mathrm{IFlag}
\]
where $\widehat{\mathrm{ZCR}}$ is the normalized zero-crossing rate, $\widehat{\mathrm{Flat}}$
is the spectral flatness, $\mathrm{CFlag}$ and $\hat{\sigma}_c$ capture spectral centroid
behavior, and $\mathrm{IFlag}$ is an insect-band autocorrelation indicator. If
$\text{Score}(f_i)$ exceeds a configurable threshold, the subframe casts a noise vote;
otherwise it is treated as bird-like. The segment label is determined by majority vote across
all six subframes:
\[
    \text{Label}(S_j) = \operatorname{mode}\bigl(\text{vote}(f_1), \ldots, \text{vote}(f_6)\bigr)
\]

\subsection{Feature Extraction (BirdNET Embedding)}

\noindent
Each bird-labeled segment is passed through the frozen BirdNET V2.4 TFLite model. The
penultimate global average pooling activation yields a fixed-dimensional embedding:
\[
    \mathbf{e}_j = \phi(S_j), \quad \mathbf{e}_j \in \mathbb{R}^{1024}
\]
where $\phi$ denotes the BirdNET feature extraction function.

\subsection{MLP Classifier and Focal Loss}

\noindent
The embedding is passed through an MLP with two hidden layers (batch normalization, ReLU,
dropout $p=0.3$) and a sigmoid output:
\[
    p = \sigma\!\left(W_2 \cdot \mathrm{ReLU}(W_1\,\mathbf{e}_j + b_1) + b_2\right)
\]
where $p$ is the probability of bird presence. Training minimizes the focal loss:
\[
    \mathcal{L} = -\alpha\,(1-p)^{\gamma}\,y\log(p)
                  - (1-\alpha)\,p^{\gamma}\,(1-y)\log(1-p)
\]
with $\gamma = 2.0$, which down-weights easy examples and focuses learning on hard cases.

\subsection{Autoencoder Reconstruction Gate}

\noindent
A symmetric autoencoder is trained on bird-only embeddings. For a given embedding $\mathbf{e}_j$,
the reconstruction error is:
\[
    \varepsilon(\mathbf{e}_j) = \frac{1}{1024}\sum_{i=1}^{1024}
                                \bigl(e_{ji} - \hat{e}_{ji}\bigr)^2
\]
where $\hat{\mathbf{e}}_j$ is the reconstructed embedding. The gating threshold
$\tau_{\mathrm{AE}}$ is set at the 95th percentile of reconstruction errors on validation
bird embeddings. Any segment with $\varepsilon > \tau_{\mathrm{AE}}$ is rejected as
out-of-distribution noise before reaching the classifier.

\subsection{Three-Band Decision Function}

\noindent
Final classification uses two fixed deployment thresholds:
\[
    \text{Output} =
    \begin{cases}
        \text{Bird},      & p \geq \tau_{\mathrm{high}} = 0.7 \\
        \text{Noise},     & p \leq \tau_{\mathrm{low}}  = 0.3 \\
        \text{Uncertain}, & \text{otherwise}
    \end{cases}
\]
Uncertain segments are routed to a human-review queue and mined for active learning.

% ---------------- 3.3 ----------------
\section{System Architecture}

\noindent
The proposed system is designed as a structured pipeline for classifying bird sounds from
environmental audio recordings. It integrates multiple stages, including preprocessing, noise
segregation, feature extraction, classification, and decision-making.

\subsection{Overall System Architecture}

\noindent
The overall architecture represents the complete flow from raw audio input to final output,
including both training and inference pipelines.

\begin{figure}[H]
\centering
\includegraphics[width=0.9\linewidth]{1.png}
\caption{End-to-End Pipeline of the Proposed Bird Sound Classification System}
\label{fig:overall_architecture}
\end{figure}

\subsection{Noise Segregation and Preprocessing}

\noindent
In the initial stage, raw audio recordings are standardized and divided into fixed-length
3-second segments. Each segment is further split into six 0.5-second subframes. Acoustic
features — zero-crossing rate, spectral flatness, spectral centroid, and insect-band
autocorrelation — are extracted from each subframe. A weighted noise score is computed and
each subframe votes bird or noise; a majority vote across subframes determines the segment
label.

\begin{figure}[H]
\centering
\includegraphics[width=0.35\linewidth]{3.png}
\caption{Noise Segregation V2: Subframe-Level Acoustic Feature Analysis and Majority Voting}
\label{fig:noise_seg}
\end{figure}

\subsection{Feature Extraction and Model Architecture}

\noindent
After noise filtering, bird-labeled segments are converted into 1024-dimensional BirdNET
embeddings. These embeddings are passed into the MLP classifier to predict the probability
of bird presence. An optional autoencoder trained on bird-only embeddings provides
reconstruction-error gating to reject out-of-distribution segments before classification.

\begin{figure}[H]
\centering
\includegraphics[width=0.7\linewidth]{4.png}
\caption{Feature Extraction, MLP Classifier, and Autoencoder Components}
\label{fig:model_arch}
\end{figure}

\subsection{Decision and Inference Flow}

\noindent
During inference, the system follows a structured decision process. If the autoencoder gate
is enabled, reconstruction error is evaluated first; segments exceeding $\tau_{\mathrm{AE}}$
are directly classified as noise. Surviving segments are scored by the MLP and routed using
$\tau_{\mathrm{low}}=0.3$ and $\tau_{\mathrm{high}}=0.7$ into confident bird, confident
noise, or uncertain bands. Uncertain segments are queued for human review and
active-learning refinement.

\begin{figure}[H]
\centering
\includegraphics[width=0.6\linewidth]{5.png}
\caption{Three-Way Decision Routing for Bird Sound Classification}
\label{fig:decision_flow}
\end{figure}

% ---------------- 3.4 ----------------
\section{Algorithms}

\begin{algorithm}[H]
\caption{Preprocessing and Noise Segregation V2}
\begin{algorithmic}[1]
\STATE Load raw audio files from the iBC53 corpus
\STATE Resample all recordings to 48~kHz mono
\STATE Segment each recording into non-overlapping 3-second clips
\FOR{each segment $S_j$}
    \STATE Compute RMS energy of $S_j$
    \IF{RMS energy $<$ $-40$~dB}
        \STATE Discard segment
    \ELSE
        \STATE Divide $S_j$ into six 0.5-second subframes
        \FOR{each subframe $f_i$}
            \STATE Extract ZCR, spectral flatness, spectral centroid (mean \& std), insect-band autocorrelation
            \STATE Compute weighted score:
                   $\mathrm{Score}(f_i) = w_{\mathrm{zcr}}\,\widehat{\mathrm{ZCR}}
                   + w_{\mathrm{flat}}\,\widehat{\mathrm{Flat}}
                   + w_{\mathrm{cflag}}\,\mathrm{CFlag}
                   + w_{\mathrm{cstd}}\,\hat{\sigma}_c
                   + w_{\mathrm{ins}}\,\mathrm{IFlag}$
            \IF{$\mathrm{Score}(f_i) >$ configurable threshold}
                \STATE Subframe votes \textit{noise}
            \ELSE
                \STATE Subframe votes \textit{bird}
            \ENDIF
        \ENDFOR
        \STATE Apply majority vote:
               $\mathrm{Label}(S_j) = \operatorname{mode}(\mathrm{vote}(f_1),\ldots,\mathrm{vote}(f_6))$
    \ENDIF
\ENDFOR
\STATE Store bird-labeled segments for embedding; retain noise-labeled segments as noise corpus
\end{algorithmic}
\end{algorithm}

\vspace{0.4cm}

\begin{algorithm}[H]
\caption{BirdNET Embedding Extraction}
\begin{algorithmic}[1]
\STATE Load frozen BirdNET V2.4 TFLite model
\FOR{each bird-labeled segment $S_j$}
    \STATE Extract penultimate global average pooling activation:
           $\mathbf{e}_j = \phi(S_j) \in \mathbb{R}^{1024}$
    \STATE Store embedding with metadata (segment origin, class label)
\ENDFOR
\STATE Build stratified train/validation/test split (75\% / 15\% / 10\%)
\end{algorithmic}
\end{algorithm}

\vspace{0.4cm}

\begin{algorithm}[H]
\caption{MLP Classifier Training with Focal Loss}
\begin{algorithmic}[1]
\STATE Load embeddings and binary labels (bird=1, noise=0)
\STATE Initialize MLP: two hidden layers, batch normalization, ReLU, dropout ($p=0.3$)
\FOR{each epoch (up to 50)}
    \STATE Forward pass; compute sigmoid output probability $p$
    \STATE Compute focal loss:
           $\mathcal{L} = -\alpha(1-p)^{\gamma}y\log p
           - (1-\alpha)p^{\gamma}(1-y)\log(1-p)$, $\gamma=2.0$
    \STATE Backpropagate and update weights (Adam optimizer, cosine-annealing schedule)
    \IF{validation F1 does not improve for 7 consecutive epochs}
        \STATE Early stopping
    \ENDIF
\ENDFOR
\STATE Sweep 50 candidate thresholds on validation set; save F1-optimal threshold
\end{algorithmic}
\end{algorithm}

\vspace{0.4cm}

\begin{algorithm}[H]
\caption{Autoencoder Gating (Bird-Only Training)}
\begin{algorithmic}[1]
\STATE Select bird-only embeddings from training split
\STATE Initialize symmetric autoencoder (encoder--bottleneck--decoder on $\mathbb{R}^{1024}$)
\FOR{each epoch (up to 30)}
    \STATE Encode embedding $\mathbf{e}_j$; reconstruct $\hat{\mathbf{e}}_j$
    \STATE Compute MSE loss:
           $\varepsilon = \frac{1}{1024}\|\mathbf{e}_j - \hat{\mathbf{e}}_j\|^2$
    \STATE Update weights
\ENDFOR
\STATE Set gating threshold $\tau_{\mathrm{AE}}$ = P95 reconstruction error on
       \textit{validation} bird embeddings
\end{algorithmic}
\end{algorithm}

\vspace{0.4cm}

\begin{algorithm}[H]
\caption{Inference with Three-Band Routing}
\begin{algorithmic}[1]
\FOR{each input segment}
    \STATE Extract embedding $\mathbf{e}_j = \phi(S_j)$
    \IF{autoencoder gate enabled}
        \STATE Compute $\varepsilon(\mathbf{e}_j)$
        \IF{$\varepsilon(\mathbf{e}_j) > \tau_{\mathrm{AE}}$}
            \STATE Route to \textbf{Noise} (OOD rejection); \textbf{skip} MLP
        \ENDIF
    \ENDIF
    \STATE Compute MLP probability $p$
    \IF{$p \geq \tau_{\mathrm{high}} = 0.7$}
        \STATE Route to \textbf{Confident Bird}
    \ELSIF{$p \leq \tau_{\mathrm{low}} = 0.3$}
        \STATE Route to \textbf{Confident Noise}
    \ELSE
        \STATE Route to \textbf{Uncertain} $\rightarrow$ human review queue
    \ENDIF
\ENDFOR
\end{algorithmic}
\end{algorithm}

\vspace{0.4cm}

\begin{algorithm}[H]
\caption{Active Learning and Iterative Refinement}
\begin{algorithmic}[1]
\STATE Collect uncertain-band segments and candidate false positives/negatives from noise output
\STATE Send curated set for expert annotation and verification
\STATE Update dataset with corrected labels
\STATE Retrain MLP (and autoencoder if needed) on augmented dataset
\STATE Repeat until performance converges
\end{algorithmic}
\end{algorithm}

% ---------------- 3.5 ----------------
\section{Summary}

\noindent
This chapter presented the complete methodology of the proposed noise-aware bird sound
classification system. The pipeline encompasses audio preprocessing with RMS-based silence
removal, the Noise Segregation V2 subframe voting mechanism, BirdNET embedding extraction,
focal-loss MLP training, autoencoder gating, and three-band decision routing with
active-learning support. The mathematical formulations align directly with the system's
implementation, and the algorithms reflect the actual thresholds and design choices used
in the system.

% ======================================================
% CHAPTER 4
% ======================================================
\newpage
\thispagestyle{plain}

\begin{flushleft}
{\fontsize{16}{18}\bfseries CHAPTER 4}
\end{flushleft}
\begin{center}
{\fontsize{16}{18}\bfseries RESULTS AND DISCUSSION}
\end{center}

\addcontentsline{toc}{section}{\textbf{CHAPTER 4}}

\vspace{0.5cm}

\setcounter{section}{0}
\renewcommand{\thesection}{4.\arabic{section}}

% ---------------- 4.1 ----------------
\section{Experimental Results}

\noindent
This section presents the performance evaluation of the proposed noise-aware bird sound
classification system. The results are compared against a baseline approach that directly uses
BirdNET detection confidence with a fixed threshold of 0.5. The evaluation is conducted on
the iBC53 dataset, with all comparisons reported on the same held-out test split of 7,193
segments, using a 75\%/15\%/10\% train/validation/test split of the full corpus.

\subsection{Performance Metrics}

\noindent
To evaluate the effectiveness of the system, standard binary classification metrics are used:
Accuracy, Precision, Recall, F1 Score, False Positive Rate (FPR -- noise misclassified as
bird), and False Negative Rate (FNR -- bird missed).

\subsection{Quantitative Results}

\noindent
Table~\ref{tab:results} summarizes the comparison between the BirdNET confidence baseline
(threshold 0.5) and the proposed noise-aware pipeline (MLP at probability threshold 0.5,
autoencoder gating omitted for parity).

\begin{table}[H]
\centering
\caption{Performance Comparison: BirdNET Baseline vs.\ Noise-Aware Pipeline on iBC53
         Test Split (N\,=\,7,193 segments)}
\label{tab:results}
\vspace{0.3cm}
\renewcommand{\arraystretch}{1.3}
\setlength{\tabcolsep}{10pt}
\resizebox{\textwidth}{!}{%
\begin{tabular}{|l|c|c|c|}
\hline
\textbf{Metric} & \textbf{BirdNET Baseline (@0.5)} & \textbf{Noise-Aware Pipeline (@0.5)} &
\textbf{Absolute Change} \\
\hline
Accuracy                                      & 0.570 & 0.992 & $+0.422$ \\
\hline
Precision                                     & 0.942 & 0.996 & $+0.054$ \\
\hline
Recall                                        & 0.558 & 0.995 & $+0.437$ \\
\hline
F1 Score                                      & 0.701 & 0.995 & $+0.294$ \\
\hline
False-Positive Rate (noise predicted as bird) & 0.315 & 0.037 & $-0.278$ \\
\hline
False-Negative Rate (bird missed)             & 0.442 & 0.005 & $-0.437$ \\
\hline
\end{tabular}%
}
\end{table}

\subsection{Performance Visualization}

\noindent
The improvement in accuracy, precision, recall, and F1 score is illustrated in
Figure~\ref{fig:performance}.

\begin{figure}[H]
\centering
\includegraphics[width=0.8\linewidth]{metrics_comparison.png}
\caption{Comparison of Accuracy, Precision, Recall, and F1 Score (N\,=\,7,193)}
\label{fig:performance}
\end{figure}

\noindent
The proposed model consistently outperforms the baseline across all evaluation metrics.

\subsection{Error Rate Analysis}

\noindent
Figure~\ref{fig:error} shows the comparison of False Positive Rate (FPR) and False Negative
Rate (FNR) between the baseline and the proposed system.

\begin{figure}[H]
\centering
\includegraphics[width=0.8\linewidth]{error_comparison.png}
\caption{Comparison of False Positive and False Negative Rates (N\,=\,7,193)}
\label{fig:error}
\end{figure}

\noindent
The false-positive rate is reduced by approximately 88\% (from 0.315 to 0.037) and the
false-negative rate by approximately 99\% (from 0.442 to 0.005), demonstrating the
effectiveness of the noise-aware approach.

\subsection{Discussion}

\noindent
The experimental results highlight the limitations of relying solely on BirdNET confidence
scores for classification. The baseline tends to misclassify noisy segments as bird sounds
and fails to detect faint bird vocalizations due to the uniform application of a single
fixed threshold. In contrast, the proposed system improves performance by integrating
Noise Segregation V2, BirdNET-based feature extraction, and a trained MLP classifier.
The use of learned embeddings allows the model to capture meaningful acoustic patterns that
a simple threshold cannot, while the structured pipeline ensures better generalization across
noisy and variable environments. The combination of preprocessing, embedding-based learning,
and three-band routing leads to a robust and reliable system for bird sound detection.

% ---------------- 4.1.6 ----------------
\subsection{System Interface (Streamlit UI)}

\noindent
A Streamlit-based graphical interface was developed to make the pipeline accessible to
non-technical users and field researchers. The interface provides end-to-end interaction
with the system, from audio upload through to classified output segments.

\subsubsection{Landing Page}

\noindent
The landing page presents the project title, a brief description of the pipeline, and
navigation controls for the different processing stages.

\begin{figure}[H]
\centering
\includegraphics[width=0.85\linewidth]{11.png}
\caption{Streamlit UI -- Landing Page}
\label{fig:ui_landing}
\end{figure}

\subsubsection{Uploading Audio}

\noindent
Users can upload one or more raw audio files in WAV or MP3 format. The interface confirms
the upload and displays basic file metadata such as duration and sample rate.

\begin{figure}[H]
\centering
\includegraphics[width=0.85\linewidth]{22.png}
\caption{Streamlit UI -- Audio Upload Interface}
\label{fig:ui_upload}
\end{figure}

\subsubsection{Audio Visualization}

\noindent
After upload, the interface renders a waveform and spectrogram of the uploaded recording,
allowing users to visually inspect the audio before processing.

\begin{figure}[H]
\centering
\includegraphics[width=0.85\linewidth]{33.png}
\caption{Streamlit UI -- Waveform and Spectrogram Visualization}
\label{fig:ui_viz}
\end{figure}

\subsubsection{Running Segmentation and Noise Segregation}

\noindent
The user initiates the preprocessing pipeline, which segments the audio into 3-second clips
and applies Noise Segregation V2. The interface displays progress and a summary of how many
segments were classified as bird-like versus noise.

\begin{figure}[H]
\centering
\includegraphics[width=0.85\linewidth]{44.png}
\caption{Streamlit UI -- Segmentation and Noise Segregation Progress}
\label{fig:ui_seg}
\end{figure}

\subsubsection{Final Segments: Bird and Noise}

\noindent
The final screen displays the classified segments organized into bird and noise categories.
Each segment can be played back directly in the browser. Confident, uncertain, and
noise-routed segments are clearly distinguished, and uncertain segments are flagged for review.

\begin{figure}[H]
\centering
\includegraphics[width=0.85\linewidth]{55.png}
\caption{Streamlit UI -- Final Classified Segments (Bird, Noise, Uncertain)}
\label{fig:ui_results}
\end{figure}

% ======================================================
% CHAPTER 5
% ======================================================
\newpage
\thispagestyle{plain}

\begin{flushleft}
{\fontsize{16}{18}\bfseries CHAPTER 5}
\end{flushleft}
\begin{center}
{\fontsize{16}{18}\bfseries CONCLUSION AND FUTURE SCOPE}
\end{center}

\addcontentsline{toc}{section}{\textbf{CHAPTER 5}}

\vspace{0.5cm}

\setcounter{section}{0}
\renewcommand{\thesection}{5.\arabic{section}}

% ---------------- 5.1 ----------------
\section{Conclusion}

\noindent
The proposed system for bird sound classification presents a structured and effective approach
to addressing the challenges associated with real-world environmental audio analysis. By
integrating preprocessing, noise segregation, feature extraction, and machine learning
techniques into a unified pipeline, the system handles complex and noisy audio data more
reliably than traditional threshold-based methods.

\vspace{0.3cm}

\noindent
A major strength of this work is the Noise Segregation V2 preprocessing step, which filters
out irrelevant and noisy segments before embedding, improving overall data quality and reducing
downstream classification errors. The use of BirdNET embeddings provides high-quality feature
representations that enable the MLP classifier to distinguish between bird sounds and noise
more effectively than raw confidence scores alone. Focal loss training further improves
performance on the imbalanced dataset by focusing learning on harder and more ambiguous examples.

\vspace{0.3cm}

\noindent
The autoencoder-based gating mechanism prevents unreliable predictions by rejecting
out-of-distribution inputs before they reach the classifier, while the three-band decision
logic ensures that genuinely uncertain cases are flagged for human review rather than forced
into an incorrect class. The active-learning component enables continuous improvement by
routing ambiguous segments to expert annotation, allowing the model to be refined iteratively
over time.

\vspace{0.3cm}

\noindent
On the iBC53 dataset test split (N\,=\,7,193 segments), the pipeline achieves accuracy of
0.992 and F1-score of 0.995, representing increases of 0.422 and 0.294 respectively over the
BirdNET confidence baseline at the same threshold, with an 88\% reduction in false-positive
rate (0.315 to 0.037) and a 99\% reduction in false-negative rate (0.442 to 0.005).
These results confirm that a structured, noise-aware pipeline substantially outperforms
single-threshold approaches in real-world acoustic monitoring scenarios.

% ---------------- 5.2 ----------------
\newpage
\section{Future Scope}

\noindent
The current system performs well in classifying bird sounds from Indian urban soundscapes,
but there is meaningful scope to extend its applicability and robustness. One important
direction is to evaluate and adapt the system across a wider variety of environmental
conditions, including recordings from different geographical regions, seasonal periods, and
recording hardware configurations. This would help assess how well the pipeline generalizes
beyond the conditions represented in the iBC53 dataset.

\vspace{0.3cm}

\noindent
Another area for improvement is the expansion and diversification of the training dataset.
Incorporating more labeled audio samples from underrepresented bird species, rare call types,
and acoustically challenging environments would improve the classifier's accuracy and reduce
its sensitivity to distribution shift.

\vspace{0.3cm}

\noindent
The system can also be enhanced by incorporating more advanced deep learning architectures,
such as transformer-based audio encoders, to capture richer temporal and spectral patterns
than the current convolutional embedding model. Additionally, the feature weights used in
the Noise Segregation V2 module are currently hand-tuned for the iBC53 corpus; making these
weights learnable from subframe-level labeled data would make the noise segregation stage
more transferable across regions and sensor types.

\vspace{0.3cm}

\noindent
Finally, the system can be extended beyond binary bird-versus-noise classification to
multi-species identification, which would significantly increase its ecological value and
support broader conservation monitoring applications.

% ======================================================
% REFERENCES
% ======================================================
\newpage
\thispagestyle{plain}

\begin{center}
{\fontsize{14}{16}\bfseries REFERENCES}
\end{center}

\addcontentsline{toc}{section}{\textbf{REFERENCES}}

\vspace{0.5cm}

%
% Citation order (by first appearance in document):
% [1]  BirdNET       — §1.1 Introduction
% [2]  WoodKahl      — §2.2 Literature Survey
% [3]  PerezGranados — §2.2 Literature Survey
% [4]  McGinn        — §2.2 Literature Survey
% [5]  Ghani         — §2.2 Literature Survey
% [6]  Navine        — §2.2 Literature Survey
% [7]  Xie           — §2.2 Literature Survey
% [8]  iBC53         — §2.2 Literature Survey
% [9]  Ntalampiras   — §2.2 Literature Survey
% [10] Kath          — §2.2 Literature Survey
% [11] Rauch         — §2.2 Literature Survey
% [12] Dumoulin      — §2.2 Literature Survey
% [13] Schwinger     — §2.2 Literature Survey
% [14] Perch2        — §2.2 Literature Survey
% [15] FocalLoss     — §2.2 Literature Survey
% [16] Hinton        — §2.2 Literature Survey
% [17] Cole          — §2.2 Literature Survey
% [18] Stowell       — §2.2 Literature Survey
%

\begin{thebibliography}{99}

\bibitem{BirdNET}
S.~Kahl, C.~M.~Wood, M.~Eibl, and H.~Klinck,
``BirdNET: A deep learning solution for avian diversity monitoring,''
\textit{Ecological Informatics}, vol.~61, Art.~no.~101236, 2021.

\bibitem{WoodKahl}
C.~M.~Wood and S.~Kahl,
``Guidelines for appropriate use of BirdNET scores and other detector outputs,''
\textit{Journal of Ornithology}, vol.~165, no.~3, pp.~777--782, 2024.

\bibitem{PerezGranados}
C.~P\'{e}rez-Granados,
``BirdNET: Applications, performance, pitfalls and future opportunities,''
\textit{Ibis}, vol.~165, no.~3, pp.~1068--1075, 2023.

\bibitem{McGinn}
K.~McGinn, S.~Kahl, M.~Z.~Peery, H.~Klinck, and C.~M.~Wood,
``Feature embeddings from the BirdNET algorithm provide insights into avian ecology,''
\textit{Ecological Informatics}, vol.~74, Art.~no.~101995, 2023.

\bibitem{Ghani}
B.~Ghani, T.~Denton, S.~Kahl, and H.~Klinck,
``Global birdsong embeddings enable superior transfer learning for bioacoustic classification,''
\textit{Scientific Reports}, vol.~13, Art.~no.~22876, 2023.

\bibitem{Navine}
A.~K.~Navine, T.~Denton, M.~J.~Weldy, and P.~J.~Hart,
``All thresholds barred: Direct estimation of call density in bioacoustic data,''
\textit{Frontiers in Bird Science}, vol.~3, Art.~no.~1380636, 2024.

\bibitem{Xie}
J.~Xie, J.~G.~Colonna, and J.~Zhang,
``Bioacoustic signal denoising: A review,''
\textit{Artificial Intelligence Review}, vol.~54, no.~5, pp.~3575--3597, 2021.

\bibitem{iBC53}
A.~Gupta, C.~S.~Swathi, A.~Jesudoss, A.~Shankar, and H.~Kumar,
``A case study: Using passive acoustic monitoring devices in a growing urban landscape
to monitor avian diversity,''
\textit{bioRxiv preprint bioRxiv:2025.07.21.665881}, 2025.

\bibitem{Ntalampiras}
S.~Ntalampiras and I.~Potamitis,
``Acoustic detection of unknown bird species and individuals,''
\textit{CAAI Transactions on Intelligence Technology}, vol.~6, no.~3, pp.~291--300, 2021.

\bibitem{Kath}
H.~Kath, P.~P.~Serafini, I.~B.~Campos, T.~S.~Gouv\^{e}a, and D.~Sonntag,
``Leveraging transfer learning and active learning for data annotation in passive acoustic
monitoring of wildlife,''
\textit{Ecological Informatics}, vol.~82, Art.~no.~102710, 2024.

\bibitem{Rauch}
L.~Rauch, D.~Huseljic, M.~Wirth, J.~Decke, B.~Sick, and C.~Scholz,
``Towards deep active learning in avian bioacoustics,''
\textit{arXiv preprint arXiv:2406.18621}, 2024.

\bibitem{Dumoulin}
V.~Dumoulin \textit{et al.},
``The search for Squawk: Agile Modeling in Bioacoustics,''
\textit{arXiv preprint arXiv:2505.03071}, 2025.

\bibitem{Schwinger}
R.~Schwinger \textit{et al.},
``Foundation models for bioacoustics -- a comparative review,''
\textit{arXiv preprint arXiv:2508.01277}, 2025.

\bibitem{Perch2}
B.~van~Merri\"{e}nboer \textit{et al.},
``Perch 2.0: The Bittern lesson for bioacoustics,''
\textit{arXiv preprint arXiv:2508.04665}, 2025.

\bibitem{FocalLoss}
T.-Y.~Lin, P.~Goyal, R.~Girshick, K.~He, and P.~Doll\'{a}r,
``Focal loss for dense object detection,''
in \textit{Proc.\ IEEE Int.\ Conf.\ Computer Vision (ICCV)}, 2017, pp.~2980--2988.

\bibitem{Hinton}
G.~E.~Hinton and R.~R.~Salakhutdinov,
``Reducing the dimensionality of data with neural networks,''
\textit{Science}, vol.~313, no.~5786, pp.~504--507, 2006.

\bibitem{Cole}
J.~S.~Cole, N.~L.~Michel, S.~A.~Emerson, and R.~B.~Siegel,
``Automated bird sound classifications of long-duration recordings produce occupancy model
outputs similar to manually annotated data,''
\textit{Ornithological Applications}, vol.~124, Art.~no.~duac003, 2022.

\bibitem{Stowell}
D.~Stowell,
``Computational bioacoustics with deep learning: A review and roadmap,''
\textit{PeerJ}, vol.~10, p.~e13152, 2022.

\end{thebibliography}

\end{document}